% !TEX root = ../../main.tex
% C3.3 — Theo vết nhiều đối tượng (RÚT GỌN 2026-09-29; bản cũ ở report/archive/)
% Nguyên lý; lý do chọn ByteTrack/BoT-SORT để ở Mục 6.1.4.
% Khớp app: config/bytetrack_tracker.json và botsort_tracker.json (cùng ngưỡng: high 0,25, low 0,1, new 0,25,
% match 0,8, fuse_score, xác nhận sau 2 lần quan sát, kết thúc khi mất dấu > 3 khung hình được xử lý hoặc > 2000 ms).
% Mốc thời gian: tệp video dùng PTS (workers/sources/file.py); RTSP dùng đồng hồ lúc worker đọc khung (workers/sources/rtsp.py)
% -> phụ thuộc tốc độ xử lý (architect §13 #7, chưa sửa).
\section{Theo vết nhiều đối tượng}
\label{sec:3.3}

\subsection{Bài toán và các thành phần cơ bản}
\label{sec:3.3.1}
\label{sec:3.3.2}

Mô hình phát hiện xử lý từng khung hình độc lập nên không biết một người có phải người đã thấy trước đó. Theo vết nhiều đối tượng (Multi-Object Tracking, MOT) gán cho mỗi đối tượng một mã giữ qua các khung hình; chuỗi khung bao của nó là một \emph{track}, trong đề tài trở thành một \emph{lần xuất hiện}, đơn vị kết quả tìm kiếm (Mục~\ref{sec:3.1.6}).

Hướng phổ biến là \emph{theo vết dựa trên phát hiện} (tracking-by-detection), với nền tảng là thuật toán SORT~\cite{bewley2016sort} gồm ba thành phần. Thứ nhất, mỗi track gắn với một bộ lọc Kalman~\cite{kalman1960} dự đoán vị trí và kích thước khung bao ở khung hình hiện tại theo giả định chuyển động gần như đều, rồi được hiệu chỉnh bằng phát hiện thực tế. Thứ hai, mức độ khớp giữa track và phát hiện được đo bằng IoU (Công thức~\ref{eq:iou}) giữa khung bao dự đoán và khung bao phát hiện. Thứ ba, việc gán mỗi phát hiện cho tối đa một track với tổng chi phí nhỏ nhất được giải bằng thuật toán Hungary~\cite{kuhn1955hungarian}; phát hiện chưa được ghép có thể khởi tạo track mới, còn track chưa được ghép được xem là đang mất dấu.

\subsection{Thuật toán ByteTrack}
\label{sec:3.3.3}

Các phương pháp trước loại bỏ phát hiện độ tin cậy thấp, nên người bị che khuất một phần dễ bị đứt track. ByteTrack~\cite{zhang2022bytetrack} liên kết hai bước: phát hiện độ tin cậy cao được ghép trước với mọi track, rồi các track còn lại ghép với phát hiện độ tin cậy thấp, nhờ đó người bị che khuất vẫn giữ track; phát hiện thấp không ghép được bị loại vì phần lớn là nền. Chỉ phát hiện độ tin cậy cao khởi tạo track, và track mất dấu được giữ thêm một lúc để nối lại.

\subsection{Thuật toán BoT-SORT}
\label{sec:3.3.4}

BoT-SORT~\cite{aharon2022botsort} mở rộng ByteTrack bằng cách điều chỉnh trạng thái của bộ lọc Kalman, bù chuyển động của camera và kết hợp đặc trưng tái nhận dạng ngoại hình (ReID) vào bước liên kết. Trong đề tài, BoT-SORT là thuật toán thay thế với cùng các ngưỡng như ByteTrack; bù chuyển động camera và ReID được tắt vì camera cố định và để giữ chi phí thấp trên máy chỉ có CPU.

\subsection{Vòng đời track trong hệ thống}
\label{sec:3.3.5}

Thuật toán theo vết nhận các phát hiện trên những khung hình đã được lấy mẫu, riêng cho từng camera, với các quy tắc sau:
\begin{itemize}
    \item \textbf{Ngưỡng độ tin cậy:} phát hiện từ 0,25 trở lên được xem là cao và có thể khởi tạo track mới; phát hiện từ 0,1 đến dưới 0,25 chỉ dùng trong bước liên kết thứ hai.
    \item \textbf{Xác nhận:} track chỉ được xác nhận sau ít nhất hai lần quan sát; track chưa xác nhận không tạo ra kết quả.
    \item \textbf{Kết thúc:} track kết thúc khi mất dấu quá ba khung hình xử lý liên tiếp hoặc quá hai giây không quan sát, hoặc khi luồng kết thúc. Cả tệp video lẫn luồng RTSP đều dùng dấu thời gian ghi trong chính luồng (Mục~\ref{sec:3.7.4}), nên quy tắc đo thời gian đã trôi qua trước camera và không phụ thuộc tốc độ xử lý của máy.
    \item \textbf{Phạm vi mã định danh:} chỉ trong một camera và một phiên xử lý; hệ thống không liên kết track của cùng một người giữa nhiều camera.
\end{itemize}
